% DELTA23 /root/five_level_six. Exact marking and decision criterion.

\subsection*{D23.F6. An explicit marked Atkin--Lehner involution}
Use the literal modular matrices
\[
 s=\begin{pmatrix}0&1\\-1&0\end{pmatrix},\quad
 r=\begin{pmatrix}0&-1\\1&-1\end{pmatrix},\quad
 \tau=sr=\begin{pmatrix}1&-1\\0&1\end{pmatrix}
\]
in $\operatorname{PSL}_2(\mathbb Z)$, and put
$H=\overline{\Gamma_0(6)/\{\pm I\}}\subset\widehat B$.
Then $H$ is the free profinite group on
\[
 a=\tau,\qquad b=r\tau^5r,\qquad
 c=r\tau^2r\tau r\tau^2r.
 \tag{D23.F6.1}
\]
Their projective matrix representatives are respectively
\[
 A=\begin{pmatrix}1&-1\\0&1\end{pmatrix},\quad
 B=\begin{pmatrix}1&-1\\6&-5\end{pmatrix},\quad
 C=\begin{pmatrix}5&-3\\12&-7\end{pmatrix}.
\]
The Atkin--Lehner matrix
\[
 W_3=\begin{pmatrix}3&-2\\6&-3\end{pmatrix},\qquad
 \det W_3=3,\quad W_3^2=-3I
\]
induces the literal involution
\[
 w(a)=c^{-1},\qquad w(b)=b^{-1},\qquad w(c)=a^{-1}
 \tag{D23.F6.2}
\]
on this marked free group. Its outer class is the $w_3$ used
in DELTA\_ON\_GT FP.17.

\begin{proof}
Here is an explicit Bass--Serre/Schreier certificate for the
free-basis assertion.  The twelve right cosets of $H$ are
the projective bottom rows modulo two and three.  Numbering
them starting with $H$ as zero, the right actions of the
generators are
\[
 s=(0\ 1)(2\ 3)(4\ 5)(6\ 8)(7\ 9)(10\ 11),
\]
\[
 r=(0\ 2\ 1)(3\ 4\ 6)(5\ 7\ 10)(8\ 11\ 9).
\]
These permutations can be checked by multiplying rows by
the displayed two matrices in $\mathbb P^1(\mathbb F_2)$
and $\mathbb P^1(\mathbb F_3)$; their stabilizer at zero
is exactly the condition that the lower-left entry vanish
modulo six.  Use two consecutive edges in each of the four
displayed $r$-cycles, starting at $0,3,5,8$, and the three
$s$-edges $(2,3),(4,5),(6,8)$ as a tree.  After collapsing
that tree and the four $r^3$ relations, precisely the three
$s$-edges $(0,1),(7,9),(10,11)$ remain.  Their oriented
Schreier loops are exactly (D23.F6.1).  The $s^2$ relations
only reverse the chosen orientation of each such edge.
Thus these loops form a discrete free basis.  Completion of
this finite-index subgroup has the induced topology from
$B$, and gives the asserted profinite free basis.

Direct multiplication yields the three displayed matrices.
It also gives, in projective matrices,
\[
 W_3 A W_3^{-1}=C^{-1},\qquad
 W_3 B W_3^{-1}=B^{-1},\qquad
 W_3 C W_3^{-1}=A^{-1}.
\]
The free-basis result now proves (D23.F6.2) as full group
equalities, not only in a finite quotient.  In particular
$W_3$ normalizes the modular subgroup.  Its shape
$\left(\begin{smallmatrix}3u&v\\6w&3z\end{smallmatrix}\right)$
and determinant three give the usual level-six Atkin--Lehner
coset for the exact divisor three.  Replacing a representative
of that coset by a member of $H$ changes the induced map only
by an inner automorphism, so the outer class is the prescribed
one.  The supplied exact program \texttt{level6\_words.py}
independently performs the row enumeration, Schreier rewriting,
matrix conjugations, and square check.
\end{proof}

\subsection*{D23.F7. Two exact word tests for the level-six comparison}
For $h\in\mathcal H$ let $\Phi_h$ be the normalized modular
automorphism fixing $\tau$, and let $\rho_h$ be E.65's
integral harmonic coordinate.  Set
\[
 \psi_h=\operatorname{Ad}(a^{8\rho_h})\Phi_h\big|_H.
 \tag{D23.F7.1}
\]
This represents $\rho_6(h)$ of FP.17 and fixes $a$.
With $y=s\tau^{-2}s$ put
\[
 k_h=\psi_h(y^{-1})y\in H,\qquad
 d_h=\operatorname{coeff}_a(k_h^{\rm ab})-
                     \operatorname{coeff}_c(k_h^{\rm ab}),
 \qquad v_h=k_hc^{d_h}.
 \tag{D23.F7.2}
\]
In this formula $\psi_h$ on $y^{-1}$ means the same conjugated
automorphism of all of $\widehat B$ before restriction.
Then the exact criterion for A.86 is
\[
 AL_{6,3}(h)=1\quad\Longleftrightarrow\quad
 \begin{cases}
 v_h w(v_h)=1,\\
 \psi_h(b)^{-1}=v_h w(\psi_h(b))v_h^{-1}.
 \end{cases}
 \tag{D23.F7.3}
\]
The balanced frame $v_h$ does not depend on which conjugator
$k_h\in H$ is used for $\psi_h(c)=k_hck_h^{-1}$.
There is no remaining unspecified inner conjugator in this test.

\begin{proof}
E.65.9--10 proves
$q_{12}\Phi_h=\operatorname{Ad}(U^{8\rho_h})q_{12}$
on the actual modular quotient, with $U=q(\tau^{-1})$.
Thus (D23.F7.1), defined first on $\widehat B$, is the identity
on its level-six quotient and preserves $H$.  It fixes $a$
since both factors do.  It represents the restriction outer
action of FP.17 because its correction is inner in
$\widehat B$.

In the literal matrices one has
\[
 g=y^{-1}=\begin{pmatrix}1&0\\2&1\end{pmatrix},
 \qquad c=ga^{-3}g^{-1}.
\]
Consequently $k_h=\psi_h(g)g^{-1}$ conjugates $c$ to
$\psi_h(c)$.  Because $\psi_h$ is the identity modulo six,
$k_h$ belongs to the principal level-six subgroup, hence to
$H$.  In terms of the original parameter this is the explicit word
\[
 k_h=a^{8\rho_h}f_h^{-1}y^{-1}f_h a^{-8\rho_h}y,
 \quad
 \psi_h(b)=a^{8\rho_h}f_h^{-1}r a^5 f_h^{-1}r a^{-8\rho_h}.
 \tag{D23.F7.4}
\]

For completeness, the following elementary free-group argument
eliminates the inner ambiguity.  Write $\psi=\psi_h$, $k=k_h$.
The condition $[\psi,w]=1$ in $\operatorname{Out}(H)$ says
that for some $v\in H$,
$\psi w=\operatorname{Ad}(v)w\psi$.
Apply this equality to $a$.  It gives
$k c^{-1}k^{-1}=v c^{-1}v^{-1}$.  Since a free generator
has centralizer its procyclic subgroup, $v=kc^e$ for a
unique $e\in\widehat{\mathbb Z}$.
Applying it to $c$ gives
\[
 k c^e w(k)\in\overline{\langle a\rangle}.
\]
If $k^{\rm ab}=p a+q b+t c$, then (D23.F6.2) shows that
the abelianization on the left is
$(p-t)a+(t+e-p)c$.
Thus necessarily $e=p-t=d_h$, and the displayed element
must be exactly $a^{d_h}$.  This is equivalent to
$v_hw(v_h)=1$.  The equality on the final free generator
$b$ is exactly the second line of (D23.F7.3).
Conversely those two equalities give the required identity
on all three free generators and hence the equality of
continuous automorphisms.  This proves the criterion.

Changing $k$ to $kc^j$ changes $d_h$ to $d_h-j$, so
$kc^j c^{d_h-j}=kc^{d_h}$.  This proves the asserted
independence.  The centralizer fact used here is the
free-factor centralizer theorem for
$H=\overline{\langle c\rangle}\amalg\widehat F_2$
(and similarly for $a$); it is the same Herfort--Ribes
Theorem B' used in D23.F1.
\end{proof}

\paragraph{What the criterion does and does not decide.}
For the element $h_*$ of E.58 one has $\rho_{h_*}=0$, so
\[
 k_{h_*}=f_{h_*}^{-1}y^{-1}f_{h_*}y,
 \qquad \psi_{h_*}(b)=f_{h_*}^{-1}r a^5 f_{h_*}^{-1}r.
\]
E.58 specifies a relative-pronilpotent lift and finite frame
coordinates, rather than these two full free-group values.
No equality or inequality in (D23.F7.3) on that element is
claimed here.  The computation supplies a marked basis and
an exact normalization for the test that was absent in FP.17;
the asserted equivalence with $\Delta=0$ remains a separate
decision.

\subsection*{D23.F8. The Atkin--Lehner matrices do not generate the division algebra of discriminant six}
An algebra generated over $\mathbb Q$ by rational
$2\times2$ Atkin--Lehner matrices is a subalgebra of
$M_2(\mathbb Q)$.  It cannot be isomorphic to the division
quaternion algebra $(2,3)$ of discriminant six.  In particular
the literal matrices used to mark A.86 do not provide an
identification with that quaternion algebra.

\begin{proof}
The Hilbert-symbol calculation $(2,3)_2=(2,3)_3=-1$,
with all other local symbols positive, is the standard
local calculation already made in D23.L3.  It is therefore
a division quaternion algebra over $\mathbb Q$.
An injective homomorphism from that four-dimensional algebra
to $M_2(\mathbb Q)$ would be an isomorphism by dimension.
But $M_2(\mathbb Q)$ has nonzero noninvertible elements and
a division algebra does not, a contradiction.  Any nonzero
unital homomorphism from a division algebra is injective,
so even that weaker possibility is excluded.
This obstruction concerns a literal matrix-algebra
identification; it does not exclude a separately constructed
Brauer class or a Shimura-curve comparison.
\end{proof}

\subsection*{D23.F10. The extra-CM element passes every pro-two level-six test}
Let $h_*$ be the actual element of LM7\_PUSH E.58, and let
$H^{(2)}$ be the maximal pro-two quotient of the marked
level-six subgroup $H$.  The automorphism $\rho_6(h_*)$
is the identity on $H^{(2)}$. Consequently each of
$AL_{6,2}(h_*),AL_{6,3}(h_*),AL_{6,6}(h_*)$ is the identity
on that entire quotient, although $\Delta(h_*)=1$.
No test factoring through the induced action on
$H^{(2)}$ can establish $AL_{6,3}=1\Longleftrightarrow\Delta=0$.

\begin{proof}
Put $U=F'F^3$, as in E.58.  That theorem supplies
$\Phi_{h_*}=1$ on $\widehat B/R_{\rm nil}(U)$ and
$\rho_{h_*}=0$.  The latter makes the correction in
(D23.F7.1) trivial.
The finite source calculation E.58.11--13, reduced modulo
six, gives
\[
 q_6(U)=\{1\}\times V_4
 \quad\hbox{in }S_3\times A_4.
\]
On the other hand $q_6(H)=C_2\times C_3$ by FP.17.
Their intersection is trivial. Thus
$V=U\cap H=\ker(q_6|_U)$ is normal of index four in $U$,
with quotient $V_4$.

For a normal subgroup $V$ of a profinite group $U$ with
finite two-group quotient, one has $O^2(V)=O^2(U)$.
Indeed, if $N$ is an open normal subgroup of $V$ of
two-power index, intersect its finitely many $U$-conjugates.
The resulting $N_0$ is normal in $U$; the quotient
$V/N_0$ embeds into a finite product of two-groups, and
$U/N_0$ is an extension of that two-group by $U/V$.
Hence $O^2(U)\subset N_0\subset N$. Intersecting over $N$
gives $O^2(U)\subset O^2(V)$.  Conversely restriction of
every finite two-group quotient of $U$ to $V$ gives
$O^2(V)\subset O^2(U)$.
Since every finite nilpotent residual lies in every
pro-prime residual, and every homomorphism from $H$ to a
finite two-group restricts to such a homomorphism on $V$,
we obtain
\[
 R_{\rm nil}(U)\subset O^2(U)=O^2(V)\subset O^2(H).
\]
For each $h\in H$ the difference
$\Phi_{h_*}(h)h^{-1}$ belongs to $R_{\rm nil}(U)$,
so its image in $H^{(2)}$ is one.  Thus the induced
automorphism is the identity.  The characteristic subgroup
$O^2(H)$ is preserved by each $w_e$, so their crossed
commutators with this automorphism are identity too.
The sign $\Delta(h_*)=1$ is E.58.23 and E.65.4.
\end{proof}

\subsection*{D23.F12. One actual mixed-sign element is invisible to every pronilpotent level-six test}
Assume the supplied relative peripheral lifting theorem
LM7\_PUSH E.58, in its stated version for every open normal
subgroup.  There is an actual element $h_6\in\mathcal H$
with
\[
 \Delta(h_6)=1,\qquad\rho_{h_6}=0,
 \qquad \rho_6(h_6)|_{H/R_{\rm nil}(H)}=1.
 \tag{D23.F12.1}
\]
Consequently all three level-six Atkin--Lehner defects of
this one element are the identity on every finite nilpotent
quotient of $H$, and on its entire maximal pronilpotent
quotient.  This is an actual gentle element, rather than
a collection of unrelated finite shadows.

\begin{proof}
Let
\[
 U=N_6=\ker(\widehat B\longrightarrow
              Q_6=\operatorname{SL}_2(\mathbb Z/6)/\{\pm I\}).
\]
Then $U$ is open normal and $U\subset H\cap F$.
Let $q=q_{12}$ be the exact projective level-twelve quotient
of E.58. Its image $q(U)=\ker(Q_{12}\to Q_6)$ is an
elementary abelian two-group of order eight.  Indeed every
class in this kernel has a unique representative congruent
to $I$ modulo six, of the form $I+6M$ with
$M\in M_2(\mathbb F_2)$ and $\operatorname{tr}M=0$.
Multiplication is addition of $M$, proving the assertion.
The central element $z_0=[5I]=[7I]$ belongs to this kernel
and corresponds to $M=I$.

Thus $q$ factors through
$\Gamma_2=\widehat B/O^2(U)$, and the kernel of
$\Gamma_2\twoheadrightarrow Q_{12}$ is pro-two.
The element $r$ has order three in both groups.
The coprime centralizer-lifting argument in the proof of
E.58.19 gives a lift
\[
 c_2\in C_{\Gamma_2}(r)\cap U^{(2)},
 \qquad q(c_2)=z_0.
\]
For specificity, that argument lifts the centralizer
element $z_0$, then conjugates its lifted conjugate of
$\langle r\rangle$ back by the pro-two kernel using
finite Schur--Zassenhaus complement conjugacy and compactness.
Since $z_0$ is identity in the retained base $Q_6$, the
corrected lift lies in $U^{(2)}$.

Use E.58.2 to identify the relative pronilpotent quotient
with the fibre product of the relative pro-prime quotients
over $Q_6$.  Give its identity automorphism the specified
peripheral conjugators
\[
 b_s=1,\qquad b_\tau=1,\qquad
 b_r=(c_2\text{ at two},\ 1\text{ at every other prime}).
\]
Their base coordinates agree, and $b_r$ really centralizes
$r$ in that quotient.  Apply the supplied relative lifting
theorem, and normalize the full lift to fix $\tau$ exactly.
This produces an actual automorphism $\Phi$ of
$\widehat B$, identity on
$\widehat B/R_{\rm nil}(U)$, with unique pure frames
\[
 \beta\in R_{\rm nil}(U),\qquad A\in U,
 \qquad q(\beta)=1,\quad q(A)=z_0.
\]
The literal parameter $f=s\Phi(s)$ lies in
$R_{\rm nil}(U)\subset U'\subset F'$.
The torsion and product calculation E.58.21, applied to
this normalized lift, gives the two gentle equations and
the actual free-group automorphism.  Hence this is an
element $h_6\in\mathcal H$.  Because $\beta\in U'\subset F'$,
its harmonic coordinate is zero.  The same signed
calculation as E.58.23 gives $(a_3,a_2)=(1,0)$ and
$\Delta(h_6)=1$.

Every continuous map from $H$ to a finite nilpotent group
restricts to such a map on $U$, since $U\subset H$.
Therefore
\[
 R_{\rm nil}(U)\subset R_{\rm nil}(H).
\]
The identity of $\Phi$ modulo the former residual proves
the last assertion in (D23.F12.1).  Each Atkin--Lehner
automorphism preserves the characteristic subgroup
$R_{\rm nil}(H)$, so its crossed commutator with this
identity action is trivial on the whole quotient.
All finite nilpotent tests follow at once from this one
full quotient identity.
\end{proof}

\paragraph{The remaining level-six gate after D23.F12.}
To distinguish this witness from the kernel of $\Delta$
using $AL_{6,3}$, a test must retain a genuinely
nonnilpotent quotient of the free level-six group.
No claim that $AL_{6,3}(h_6)$ itself is trivial is made.
